\begin{table}[t]
\caption{Leakage audit performed on the released splits. Each check is
run on the corpus and its outcome reported, rather than asserted. The
first four checks are exact and all pass. The fifth, near-duplicate
detection by perceptual hashing, is reported as a \emph{diagnostic}
rather than a gate: on a corpus in which every document shares one
of six layout templates, the same-document and different-person
distance distributions overlap, so no threshold can separate a
leaked duplicate from two unrelated documents. See
Section~\ref{subsec:leakage} and Fig.~\ref{fig:phash}.}
\label{tab:leakage}
\centering
\scriptsize
\setlength{\tabcolsep}{3.5pt}
\begin{tabular}{@{}p{1.15cm}p{3.0cm}ll@{}}
\toprule
Leakage mode & Check & Required & Observed \\
\midrule
Identity & person-id intersection across splits & empty & 0 \\
Identifier & Aadhaar/PAN string intersection & empty & 0 \\
Template & every layout variant in every split & all & 6/6 \\
Near-duplicate & \dsPhashBits-bit perceptual-hash distance & (diagnostic) & see Sec.~\ref{subsec:leakage} \\
Augment\-ation & degradation draws of one document & one split & 0 span splits \\
Printed values & identical field-value set in two splits & none & 0 of 2100 \\
\bottomrule
\end{tabular}
\end{table}
